%=======================================================================
%  limitative_draft.tex  —  NEW backbone for the limitative paper
%  Drafted 2026-07-07.  Slots into main.tex; overwrites nothing.
%
%  Contents:
%    (A) revised abstract  (replaces the current abstract, L16–33)
%    (B) Lemma (inertness of the alignment operator)  — insert in §4
%    (C) Theorem 3'  (the functional no-go)  — REPLACES the broken §5–6 Thm 3
%    (D) Corollary   (what external input is required)  — feeds the outlook §7
%
%  Preamble needed (already in main.tex or add): amsmath, amssymb, amsthm,
%  with \newtheorem{theorem}, {lemma}, {corollary} and the proof environment.
%  Notation: E = the 2D degenerate mu-tau eigenspace of H_nu; H_e=Y_eY_e^dag,
%  H_nu=m_nu m_nu^dag; 1_E the identity on E; R(theta) the residual SO(2).
%=======================================================================


%========================= (A) REVISED ABSTRACT ========================
\begin{abstract}
The finite spectral triple of the Standard Model geometrizes fermion mass
spectra but does not, in the presence of degeneracies, fix the relative
orientation of generation subspaces. We show that whenever the effective
neutrino Hermitian operator $H_\nu = m_\nu m_\nu^\dagger$ admits a
two-dimensional eigenspace $E$, an associated residual $\mathrm{SO}(2)$ freedom
in the $\mu$--$\tau$ plane is invisible to every axiom of noncommutative
geometry and to the spectral action, yet acts nontrivially on the PMNS matrix
and carries the \emph{entire} indeterminacy of the atmospheric mixing angle
$\theta_{23}$ (Theorem~\ref{thm:obstruction}). Subject to the order-zero,
first-order, and reality conditions, the unique minimal relational operator
built linearly from the sector Hermitians and capable of lifting this freedom
is $A(\alpha,\beta)=\alpha H_e+\beta H_\nu$
(Theorem~\ref{thm:minimal}). We show, however, that on the degenerate subspace
this operator is \emph{inert}---its $H_\nu$ component acts as a multiple of the
identity---so its eigenbasis coincides with that of $H_e$ and selects the
\emph{aligned} configuration $\theta_{23}=0$
(Lemma~\ref{lem:inertness}). More strongly, we prove that \emph{no} functional
of the misalignment (commutator-norm) class can select the physical value: for
any such functional the stationary configurations are rigidly
$\{0,\pi/4,\pi/2\}$, independent of all fermion mass magnitudes, with the
minima at exact alignment and the value $\pi/4$ a symmetry-forced maximum
rather than a derived quantity (Theorem~\ref{thm:nogo}). We conclude that
$\theta_{23}$ is not determined by spectral geometry in the (near-)degenerate
regime; it is fixed only by $\mu$--$\tau$--symmetry--breaking data external to
the finite triple, and we characterize precisely the structure such input must
supply (Corollary~\ref{cor:external}). These results delineate the boundary of
what spectral geometry alone can fix in the lepton flavor sector.
\end{abstract}


%============ (B) LEMMA — inertness of the alignment operator ===========
% Insert in Section 4, immediately after A(alpha,beta)=alpha H_e+beta H_nu
% is introduced and shown to be the minimal admissible operator.

\begin{lemma}[Inertness of the alignment operator on $E$]
\label{lem:inertness}
Let $E\subset\mathbb{C}^3$ be the two-dimensional eigenspace on which
$H_\nu|_E=\lambda\,\mathbf 1_E$, and let $H_e|_E$ be nondegenerate
($m_\mu\neq m_\tau$). Then for every $\alpha\neq0$ and every $\beta\in\mathbb R$,
\[
  A|_E \;=\; \alpha\,H_e|_E \;+\; \beta\lambda\,\mathbf 1_E ,
\]
so $A|_E$ and $H_e|_E$ have identical eigenprojectors. Consequently the
orientation of $E$ selected by diagonalizing $A$ is exactly that of $H_e$, and
adopting it as the neutrino frame gives a trivial $\mu$--$\tau$ block of the
PMNS matrix, i.e. $\theta_{23}=0$.
\end{lemma}

\begin{proof}
By definition of $E$, $H_\nu|_E=\lambda\,\mathbf 1_E$, hence
$A|_E=\alpha H_e|_E+\beta\lambda\,\mathbf 1_E$. Adding a scalar multiple of the
identity shifts eigenvalues but leaves eigenvectors unchanged, so for
$\alpha\neq0$ the spectral projectors of $A|_E$ coincide with those of
$H_e|_E$. Writing $U_{\mathrm{PMNS}}=U_e^\dagger U_\nu$ and choosing $U_\nu$ on
$E$ to be the eigenframe of $A|_E$ (equivalently of $H_e|_E$) gives
$U_e^\dagger U_\nu|_E=\mathbf 1_E$ on the $\mu$--$\tau$ block, i.e.
$\theta_{23}=0$.
\end{proof}

\noindent\emph{Remark.} The $H_\nu$ term is thus dynamically silent precisely
where the residual freedom lives: on $E$ the ``relational'' operator collapses
to $H_e$ alone. This is the structural reason a minimization of misalignment
returns the unmixed configuration, and it motivates the sharper statement
below.


%============ (C) THEOREM 3' — the functional no-go =====================
% Replaces the previous Section 5-6 "Theorem 3" (unique minimizer -> theta23).
% Set delta_e = m_tau^2 - m_mu^2 > 0, and delta_nu >= 0 the mu-tau splitting of
% H_nu on E (delta_nu = 0 is the exact-degeneracy limit of Theorem 1).
% psi_0 is the relative eigenbasis angle of H_e and H_nu on E; in the
% alignment convention theta_23 = psi_0 + theta.

\begin{theorem}[No misalignment functional determines $\theta_{23}$]
\label{thm:nogo}
Work on $E\cong\mathbb{C}^2$. Let $H_e|_E$ have eigenvalue gap
$\delta_e=m_\tau^2-m_\mu^2>0$ and let $H_\nu|_E$ have gap $\delta_\nu\ge0$, with
relative eigenbasis angle $\psi_0$. Let $R(\theta)$ be the residual rotation of
$E$ and define the misalignment functional
\[
  \mathcal{E}(\theta)\;=\;\bigl\|\,[\,H_e,\;R(\theta)^\dagger H_\nu R(\theta)\,]\,\bigr\|_F^2 .
\]
Then:
\begin{enumerate}
\item[(i)] If $\delta_\nu=0$ (exact degeneracy, the hypothesis of
Theorem~\ref{thm:obstruction}), $[\,H_e,R(\theta)^\dagger H_\nu R(\theta)\,]=0$
for all $\theta$; $\mathcal E\equiv0$ possesses no isolated extremum and selects
no orientation.
\item[(ii)] If $\delta_\nu>0$,
\[
  \boxed{\;\mathcal{E}(\theta)\;=\;C\,\sin^2\!\bigl(2(\theta+\psi_0)\bigr),
  \qquad C=\tfrac12\,\delta_e^2\,\delta_\nu^2\;>\;0. \;}
\]
Its stationary points are $\theta+\psi_0\in\tfrac{\pi}{4}\mathbb{Z}$; reading
$\theta_{23}=\psi_0+\theta$, the accessible stationary values are exactly
\[
  \theta_{23}\in\{\,0,\ \tfrac{\pi}{4},\ \tfrac{\pi}{2}\,\}\pmod{\tfrac{\pi}{2}},
\]
the minima ($\mathcal E=0$) being the aligned configurations $\theta_{23}=0,\pi/2$
and $\theta_{23}=\pi/4$ the unique interior maximum.
\item[(iii)] The amplitude $C$ carries the entire dependence on the fermion
mass magnitudes and is a strictly positive prefactor; it does not enter the
location of the stationary points. Hence for any monotone $F$, every extremizer
of $F(\mathcal E)$ lies in $\{0,\pi/4,\pi/2\}$, independent of all mass inputs.
\end{enumerate}
\end{theorem}

\begin{proof}
Decompose each restricted operator into trace and traceless parts. A $2\times2$
Hermitian $X$ with eigenvalue gap $\delta_X$ may be written
$X=c_X\mathbf 1+\tfrac{\delta_X}{2}\,\hat n_X\!\cdot\!\vec\sigma$, with
$\hat n_X$ a unit Bloch vector in the plane spanned by $\sigma_x,\sigma_z$ at
polar angle $2\xi_X$, where $\xi_X$ is the eigenbasis orientation of $X$ in $E$.
Trace parts and the scalar $c_X$ commute with everything, so they drop from the
commutator:
\[
  [H_e,H_\nu]=\Bigl[\tfrac{\delta_e}{2}\hat n_e\!\cdot\!\vec\sigma,\
  \tfrac{\delta_\nu}{2}\hat n_\nu\!\cdot\!\vec\sigma\Bigr]
  =\frac{\delta_e\delta_\nu}{4}\,[\hat n_e\!\cdot\!\vec\sigma,\ \hat n_\nu\!\cdot\!\vec\sigma]
  = i\,\frac{\delta_e\delta_\nu}{2}\,(\hat n_e\times\hat n_\nu)\!\cdot\!\vec\sigma,
\]
using $[\,\hat a\!\cdot\!\vec\sigma,\hat b\!\cdot\!\vec\sigma\,]=2i(\hat a\times\hat b)\!\cdot\!\vec\sigma$.
With $\operatorname{tr}(\sigma_i\sigma_j)=2\delta_{ij}$,
\[
  \|[H_e,H_\nu]\|_F^2=\frac{\delta_e^2\delta_\nu^2}{4}\cdot 2\,|\hat n_e\times\hat n_\nu|^2
  =\frac{\delta_e^2\delta_\nu^2}{2}\,\sin^2\!\bigl(2(\xi_e-\xi_\nu)\bigr),
\]
since $\hat n_e,\hat n_\nu$ lie in a common plane at angles $2\xi_e,2\xi_\nu$, so
$|\hat n_e\times\hat n_\nu|=|\sin(2\xi_e-2\xi_\nu)|$. Writing the relative angle
$\psi_0=\xi_e-\xi_\nu$ gives the exact-degeneracy statement (i): if
$\delta_\nu=0$ the commutator vanishes identically. For $\delta_\nu>0$, the
residual rotation acts by $\xi_\nu\mapsto\xi_\nu-\theta$, i.e.
$\psi_0\mapsto\psi_0+\theta$, and adding it back does not affect the trace parts;
hence
\[
  \mathcal E(\theta)=\tfrac12\delta_e^2\delta_\nu^2\,\sin^2\!\bigl(2(\theta+\psi_0)\bigr),
\]
which is (ii) with $C=\tfrac12\delta_e^2\delta_\nu^2$. Stationarity:
$\mathcal E'(\theta)=2C\sin\!\bigl(4(\theta+\psi_0)\bigr)=0\iff\theta+\psi_0\in\tfrac{\pi}{4}\mathbb Z$;
$\mathcal E=0$ (minima) at $\theta+\psi_0\in\tfrac{\pi}{2}\mathbb Z$ and the
maximum at the half-integer points, giving the stated $\theta_{23}$ values.
Finally $C>0$ multiplies $\mathcal E$ uniformly, so $\operatorname*{arg\,ext}_\theta F(\mathcal E(\theta))$
is unchanged by $C$ for any strictly monotone $F$, which is (iii).
\end{proof}

\begin{corollary}[External input required to fix $\theta_{23}$]
\label{cor:external}
No functional built monotonically from the misalignment
$\|[H_e,R(\theta)^\dagger H_\nu R(\theta)]\|$ selects a value of $\theta_{23}$ in
$(0,\pi/2)\setminus\{\pi/4\}$; in particular it cannot select the observed
$\theta_{23}\approx49^\circ$. Minimization returns exact alignment,
$\theta_{23}=0$ (Lemma~\ref{lem:inertness}); the value $\pi/4$ is the
$\mu$--$\tau$ (2$\leftrightarrow$3 exchange) symmetric point, a symmetry-forced
maximum independent of the mass inputs, not a determination. To select an
interior value one must augment the construction with a Hermitian operator $Q$
on $E$ such that
\begin{enumerate}
\item[(a)] $Q|_E\not\propto\mathbf 1_E$ (so $Q$ is \emph{non-inert}, unlike
$H_\nu$ on $E$), and
\item[(b)] the eigen-orientation of $Q$ on $E$ is fixed by data other than
$\theta_{23}$ itself,
\end{enumerate}
for instance a controlled $\mu$--$\tau$--breaking of the neutrino sector whose
orientation is pinned by the mass-squared splittings $\Delta m^2_{21},\Delta
m^2_{31}$ or by renormalization-group running between the seesaw and electroweak
scales. Absent such an operator, $\theta_{23}$ is a datum of the theory, not a
prediction of it.
\end{corollary}

\noindent\emph{Remark.} The dichotomy is sharp: a \emph{symmetry}
(2$\leftrightarrow$3 exchange) forces exactly $\pi/4$, while a
\emph{relational functional} of the above class has its stationary set frozen at
$\{0,\pi/4,\pi/2\}$. The physical deviation $\theta_{23}-\pi/4$ is precisely the
$\mu$--$\tau$--breaking that neither device supplies; it is the content of
condition (b) above.
